% !TeX root = incremental_SS_Translation.tex
% Paras

\chapter{Cikitsāsthāna 5:  On the Treatment of Great Wind Diseases}

\section{Literature}

Meulenbeld offered an annotated overview of this chapter and a bibliography
of earlier scholarship to 2002.\fvolcite{IA}[266]{meul-hist}

\section{Translation}

\begin{translation}

\subsection{Wind-afflicted blood}
\item [1]
Now we shall describe the treatment of great wind diseases.



\item [3]

Some people say that there are two kinds of
\se{vātarakta}{wind-afflicted blood}, the
\se{uttāna}{superficial} and \se{avagāḍha}{deep}.\footnote{On
\se{vātarakta}{wind-afflicted blood}, see the discussion at
\pageref{ni1:vatarakta}.

Ḍalhaṇa commented \citep[424]{vulgate} that \dev{uttāna} refers to
being situated in the skin and flesh, and \dev{avagāḍha} refers to
being situated internally.

The view referred to can be found at \CS\ \Ca{1.19.3}{109--110} and
\Ca{6.29.19}{628}.  See the contextual discussion by
\citet[101]{emme-1984}.} However, this is not correct.

Why?

It is like \se{kuṣṭha}{pallid skin disease}: having been superficial,
it becomes deep after an interval of time.  Therefore counting it as
two things is rejected.

\item[4]

In this regard, it brings about pains that are caused by
\se{vātaśoṇita}{wind-afflicted blood}.

The wind of someone who habitually eats too much, and heavy, hot food,
gets irritated by specific things like conflicts with  strong
people.\footnote{\Dalhana{4.5.4}{424} noted that “conflicts with
    powerful people, etc." were listed in the \emph{Vraṇapraśna} chapter
    (\SS\ 1.21); see \SS\ \Su{1.21.19}{103}. A long list of wind-irritants
    is given in that passage.  \Dalhana{1.15.32}{73}, \emph{et passim},
    glossed \dev{adhyaśana} not as “overeating” but as “eating when
    previous food is not yet digested."} The wind's path is obstructed by
    corrupted blood that suddently unites with it and causes pains that
    are triggered by blood with wind in it.  Experts refer to that as
    “\se{vātaśonita}{wind-afflicted blood}.”  It initially establishes
    itself in the hands and feet.  Then it spreads throughout the body.
    Its preliminary symptoms are pricking, burning, itching, swelling,
    stiffness, rough skin, throbbing of the ducts, sinews and pipes, and
    weakness of the thigh. There is also the chance appearence of brown
    and red \se{maṇḍala}{round patches} that arise on the soles of the
    feet, fingers, ankles and wrists. This disease fully manifests in
    people who do not have treatment and who do not behave as they
    should.\footnote{This specific type of “treatment” (\emph{pratisāra})
        is described as being of four types at \SS\ \Su{4.40.69}{558}, i.e.,
        paste, linctus, honey and powder \pvolcite{1}[527]{josi-maha}.}

Its characteristics have been stated.   In the case of any of these,
those who do not take precautions develop disabilty.


\item[5]

Wind-blood most commonly gets irritated in those who are very
delicate, those who have wrong diet and activity, those who are
corpulent,
and those who live in comfort.\footnote{This same passage
occurs in \Su{2.1.40}{263}, where Ḍalhaṇa questioned its
appropriateness given its repeated occurrence here.  See footnotes
\ref{ni1:vatarakta} and \ref{2.1.40} (pages
\pageref{ni1:vatarakta}
and \pageref{2.1.40}).}

\item[6]

In such a case, one should treat a patient if he is strong,
self-possessed, and well-resourced, as long as he is not being
\se{ava\root śṝ}{worn down} by breathlessness, thirst, fever,
fainting, dyspnea, stiffness, or \se{arocaka}{loss of interest in
food},  and not secondarily afflicted by
\sepl{saṅkoca}{cramp}.\footnote{On \sepl{upadrava}{secondary
affliction}, see a floating definition verse cited by
\Dalhana{6.39.296ab}{693} and the definition in \CS\
\Ca{6.21.40}{561}. The difficulties in rendering this term were
discussed by \citet[447--448]{meul-1974}.

The syntax of the dvandva compound ending in the present passive
participle \dev{anavaśīryamāṇam} is unusual; it is written out of the
vulgate.}

\subsubsection{Therapy}
\item[7]

In this case, for fear of irritating the wind, one should start by
repeatedly draining the corrupted blood, little by little, that is
blocking the \sepl{mārga}{channel},  except when the person's body is
weakened by excessive wind and dryness.\footnote{“Drain” translates
    \dev{avasiñcet}, which is glossed by \Dalhana{4.5.7}{425} as
    “\dev{āsravet}.” See also \volcite{1}[64]{josi-maha}, which also
    glosses \dev{avasecana} as \dev{visrāvaṇam}.  Bloodletting is being
    indicated.} Then, one should apply the preliminary treatments like
    emesis.\footnote{The following therapies are for different types of \se{vātaśonita}{wind-afflicted blood} according to humoral preponderance.}
    
\paragraph{Strong wind}    
    Next, \diff{if the wind, mixed with food, is strong}, the
    patient should be given mature ghee to drink.

Alternatively, he should be made to drink goat's milk, with half its
quantity of oil, mixed with \se{akṣa}{twelve grams} of
\gls{madhuka},\footnote{\Dalhana{4.5.7}{425} clarified that “half the
    quantity” meant half the total quantity of the milk. For the measure
    \dev{akṣa}, which is a synonym for \dev{karṣa}, see
    \cite[263]{wuja-2003}.}  or else it can be cooked with
    \gls{śṛgālavinnā} and sweetened with sugar and honey. Or it could be
    cooked with \gls{śuṇṭhī} and \gls{kaśeru}.

Alternatively, one should cook sesame oil together with and an
\se{pratīvāpa}{admixture} of a \se{kalka}{paste} of \gls{meṣaśṛṅgī},
\gls{śvadaṃṣṭrā}, honey,\footnote{The vulgate reads \egls{madhūka} for
    “honey."} \gls{nāgabalā}, \gls{bhadradāru}, \gls{vacā}, and
    \gls{surabhi}, together with milk with a decoction that has been
    processed eight times, consisting of \gls{śyāmā}, \gls{rāsnā},
    \gls{suṣavī}, \gls{śṛgālavinnā}, \gls{pīlu}, \gls{śatāvarī},
    \gls{śvadaṃṣṭrā}, and the \gls{dvipañcamūlaka}, and use this in drinks
    and the like.\footnote{The vulgate version of this passage seems to
        contain dittography of \egls{daśamūla} and is slightly incoherent
        (\cite[424]{vulgate}).}

Alternatively, an admixture of \gls{kākolī}, etc., is
cooked with a decoction of
\gls{śatāvarī}, \gls{apāmārga},
%\footnote{Ḍalhaṇa glossed \citep[425]{vulgate} 
%\textit{mayūraka} as 
%\textit{apāmārga}.},
\gls{dhavaka}, \gls{madhuka}, \gls{kṣīravidārī},
\gls{balā},
\gls{atibalā}, and \gls{tṛṇapañcamūlī}.\footnote{This
recipe is different in
the vulgate.  The “admixture of \gls{kākolī} etc.” is
an item that occurs
elsewhere in other recipies, e.g.,
\Su{6.17.94}{632}.}
%    \footnote{Ḍalhaṇa commented
%        \citep[425]{vulgate} that \gls{kuśa}, 
%\gls{kāśa}, \gls{nala},
%        \gls{darbha}, \gls{kāṇḍa}, and \gls{ikṣuka} 
%are called \textit{tṛna}
%        (grass).}, 


Alternatively, one may use the \gls{balā}-oil that is prepared by the
\se{śatapāka}{hundred-cook} method.\footnote{The “hundred-cook” method
    is something cooked a hundred times or sometimes with a hundred
    ingredients. See p.\,\pageref{satapaka} above. See also
    \volcite{1}[803]{josi-maha}.}

That said, one should \se{pariṣecana}{affuse} the patient with sours
or with milk that has been cooked with roots that remove
wind.\footnote{Ḍalhaṇa glossed “\se{amla}{sours}” as “items like
    liquor, \se{sauvīraka}{sour gruel}, \se{tuṣodaka}{rice-gruel}, etc”
    (\dev{surāsauvīrakatuṣodakādibhiḥ} \Dalhana{4.5.7}{425}).  The text
    does not make clear what it is that is being moistened.

We have added “the patient” as the recipient but the syntax would
allow the object to be the site of the blood letting.}

%    , which referse
%    to \CS\ \Ca{3.6.17}{256}.}

%\emph{Śatapāka} seems to be an oil that is
%prepared with a hundred parts of some things 
%similar to
%\textit{sahasrapāka} that is prepared with 
%one thousand parts of some
%herbs. Refer \textit{Cikitsāsthāna} Ch. 4 text 
%29 for the preparation
%of \textit{sahasrapāka}.}
% Draft:



%Alternatively, [the affected body part] should 
%be moistened with milk
%that is boiled with the roots of 
%wind-alleviating herbs, or it should
%be moistened with sour things. 
%%%\footnote{Ḍalhaṇa commented
%    \citep[425]{vulgate} that the sour things 
%(\textit{amla)} are
%    \gls{surā}, \gls{sauvīraka}, 
%\gls{tuṣa}-water, etc. \textit{Surā}
% is
%    some kind of liquor, \textit{sauvīraka} is 
%perhaps the fruit of
% the
%    jujube tree, and \textit{tuṣa} is perhaps 
%Terminalia Bellerica
%    (\dev{vibhītaka}).}
In that regard,
when
barley,
wheat, sesame,
\gls{mudga} and
\gls{māṣa} are ground and mixed individually
with a paste made of
%
\gls{kākolī},
\gls{kṣīrakākolī},
\gls{jīvaka},
\gls{ṛṣabhaka},
\gls{balā},
\gls{atibalā},
\gls{pṛśniparṇī},
%\footnote{\emph{śṛgālavinnā}},
\gls{meṣaśṛṅgī},
\gls{piyāla},
sugar,
\gls{kaśerukā},
\gls{surabhi}, and
\gls{vacā}
in order to make a \se{upanāha}{poultice},
then the five, which are cooked in ghee, oil, fat,
marrow,
and milk, are described as milk-based items.

%Alternatively, a \se{veśavāra}{minced 
%meat}\footnote{\volcite{1}[787]{josi-maha}.}
%    ground with the juice of oily fruit, 
%    a poultice,
%    fire,
%    and fatty fish.



Alternatively, a \se{vaiśavāra}{spiced mince}
made from ground \gls{nalamīna}
%
%\footnote{H has
%        the compound word 
%\dev{nalapīnamatsya}. \dev{nalamīna} is a
%particular
%        fish known as \textit{cilicima} 
%(\dev{cilicimaḥ}). See
%        \textit{Amarakośa}. Also, if the name is 
%\dev{nalamatsya}
% then the
%        word \dev{pīna} (fat) within the name is 
%not according to
% proper
%        Sanskrit. But, it can be allowed because 
%the word
% \dev{matsya} (fish),
%        instead of being a part of the name, can 
%be considered to
% mean fish in
%        general and thus the word \dev{pīna} 
%becomes its modifier.
% Thus,
%        \dev{nalapīnamatsya} can mean "a fat 
%fish that is a
% \dev{nala}
%        (\textit{cilicima})".\\ Ḍalhaṇa says in his 
%comment
%        \citep[425]{vulgate} that 
%\dev{nalamīna} is a type of
% \dev{rohita}
%        (\textit{rohita}). Monier Williams says 
%that \textit{rohita}
% is a kind
%        of fish: Cyprinus Rohitaka. Regarding 
%the \textit{rohita}
% fish, there
%        is a \textit{subhāṣita}: 
%\dev{agādhajalasañcārī na garvaṃ
% yāti rohitaḥ
%        | aṅguṣṭhodakamātreṇa śapharī 
%pharpharāyate ||} This
% indicates that
%        \textit{rohita} is a deep water 
%fish.}\q{The webpage
%
%https://hindi.shabd.in/vairagya-shatakam-bhag-acharya-arjun-tiwari/post/117629
%             says that this verse belongs to the 
%\textit{Nītiratna}.
% I could not find
%            this text.
can be used instead.\footnote{On \dev{vaiśavāra} “minced meat” see
    \volcite{1}[787]{josi-maha}, citing \CS\ \Ca{1.27.269}{169} and \SS\
    \Su{4.2.25}{409}.  The Nepalese manuscripts use the form
    \dev{vaiśavāra} (also \CS\ \Ca{6.25.73}{595}) while other sources and
    dictionaries cite \dev{veśavāra, veṣavāra} or \dev{vesavāra}. The term
    has received much discussion amongst the commentators. Cakrapāṇidatta,
    on \Ca{1.27.268--270}{169}, described \dev{veśavāra} by citing a
    treatise on cookery:  “According to tradition, in the discipline of
    Cookery, the proper \dev{veśavāra} is: one should cook deboned meat,
    well-steamed and repeatedly ground on a stone, and mixed with long
    pepper, ginger, black pepper, sugar, and ghee all together
    (\dev{veśavāraḥ sūdaśāstre -- “māṃsaṃ nirasthi susvinnaṃ punar dṛṣadi
    peṣitam/ pippalīśuṇṭhimaricaguḍasarpiḥsamanvitam// aikadhyaṃ vipacet
    samyag veśavāra iti smṛtaḥ/”}).  The vulgate text of the \SS,
    \Su{1.46.365--366ab}{240}, has this exact passage, although it is not
    present in the Nepalese version. Cakrapāṇidatta did not say he was
    citing the passage from the \SS, so it may not have been in the
    version of the text available to him.  Ḍalhaṇa (\emph{idem}) glossed
    \dev{veśavāra} in much the same terms, as, “deboned meat that has been
    ground, oiled, and steamed” (\dev{nirasthisvinnasnigdhapiṣṭamāṃsa}).
    \citet[205]{meul-1974} described a variant of this dish based on the
    the \emph{Madhukośa} and gave the term a detailed discussion at
    \emph{ibid.}, pp.\,502--504.

Witness H reads the hapax legomenon
\dev{nalapīna} as the name of the fish in
this dish, but \Dalhana{4.5.7}{425}
reported reading \dev{nalamīna} at this
point in some manuscripts, identifying it as a
type of carp, see \egls{nalamīna}, and we
adopt that attested reading.}

\q{Cite Hoernle 1893-1912 (II)
35.177.<Hoernle 1908: 325>}



Alternatively, [one can use] the poultice
containing
\gls{bilva}-rind,
%\footnote{The word \dev{pesikā} in H 
%should be read \dev{peśikā}.}, 
\gls{tagara},
\gls{devadāru},
\gls{saralā},
\gls{rāsnā},
\gls{hareṇu},
\gls{kuṣṭha},
\gls{śatapuṣpa},
liquor, yogurt, and whey.

Alternatively,  [one can use] the ointment prepared by mixing
\gls{mātuluṅga}, \emph{amla}, %    \footnote{Perhaps it could mean
% vinegar
%or
%        sour curds. Refer to Monier Williams
%Sanskrit Dictionary.},
salt, and ghee with honey and \gls{śigru}-root. Or else, [one can use]
the unctuous sesame paste. 


\paragraph{Strong bile}    

\item[8]

When the [condition of wind-blood] has a predominance of bile, the
patient should be made to drink a decoction of grapes, \gls{ārevata},
\gls{kaṭphala}, \gls{payasyā}, \gls{madhuka}, \gls{candana}, and
\gls{kāśmarī}. This decoction is sweetened with honey and sugar before
consumption. Or, the decoction of \gls{śatāvarī}, \gls{paṭola},
\gls{patra}, \gls{triphalā}, \gls{kaḍurohiṇī}, and \gls{guḍūcī} should
be given. [The patient should be administered] ghee that is prepared
with sweet, bitter, and astringent [remedies].\footnote{Ḍalhaṇa
    commented \citep[425]{vulgate} that the sweet remedies are
    \gls{kākolī}, etc., bitter remedies are \gls{paṭola}, etc., and
    astringent remedies are \textit{triphalā}, etc.}

[The patient] should be sprinkled with a decoction of
\gls{bisamṛṇāla}, \gls{bhadraśriya}, and \gls{padmaka} mixed with
goat-milk\footnote{The compound word ending with \dev{kaṣāyeṇa} is
    taken to be a \textit{bahuvrīhi} for \dev{ajākṣīreṇa} (goat-milk).},
    or with rice water that is mixed with milk, sugarcane juice, honey,
    and sugar, or with whey and sour rice gruel mixed with a decoction of
    grapes and sugarcane. Or else, [the patient] should be sprinkled with
    fermented gruel mixed with whey and with ghee that is prepared with
    \gls{jīvanīya} or sprinkled with ghee that is purified for one
    hundred times.\footnote{At \Su{1.45.214-216}{213}, Ḍalhaṇa glossed
        \emph{dhānyāmla} (``fermented gruel'') as a type of gruel
        (\emph{kāñjika}) that is alcohol made from ingredients such as rice,
        \gls{jūrṇāhva}, and \gls{kodrava} (\dev{dhānyāmlaṃ kāñjikaṃ
        śālijūrṇāhvakodravādikṛtaṃ madyam}). The term \emph{jīvanīya} is a
        group of medicinal herbs, which are listed in various works, including
        \CS\ \Ca{1.4.9.1}{32}. There is also an Ayurvedic preparation called
        \emph{jīvanīya-ghṛta}; see \volcite{1}[349]{josi-maha}.\label{jīvanīya}} % reconsider
        % punctuation and how to read śatadhautena or consider emending to the
        % vulgate śatadhautaghṛtena.

Alternatively, a poultice may
be made of ground rice flour or mashed fermented rice gruel mixed with
\gls{nala}, \gls{vañjula},
\gls{tālīsa},
\gls{śṛṅgāṭaka}, 
\gls{kaloḍya},
\gls{gaurī}, \gls{śaivāla},
\gls{padma}, etc.\footnote{On fermented  gruel see \cite[23, et passim]{mchu-2021}.

The reading \emph{kālodyā} of witness H is not elsewhere attested. We
have emended to the reading of the \citep[425]{vulgate}. It is
possible, however, that \emph{kālodyā} was a local variant in use in
fifteenth-century Nepal.} The poultice should be mixed with ghee.


\paragraph{Strong blood}    

\item[9]

It is just the same in the case of a predominance of blood. And
one should perdorm blood-letting repeatedly.


\paragraph{Strong phlegm}  

\item[10]

But, in the case of a predominance of
phlegm, one should make the patient 
consume a decoction of
\gls{āmalaka} and \gls{haridrā} that is sweetened with honey.  

Alternatively, a decoction of \textit{triphalā}.

Alternatively, a paste of \gls{madhuka}, \gls{śṛṅgavera},
\gls{harītakī}, and \gls{tiktarohiṇī} with honey. 

Alternatively, he should be made to drink
\gls{harītakī} either with water or urine.   He should be
affused with oil, urine, \se{kṣārodaka}{alkaline water},  \se{surā}{liquor}, and \se{śukta}{sour gruel}. 
   
Alternatively, he should be affused  with a decoction of \gls{āragvadha}, etc.

Alternatively he should be massaged with 
with a ghee cooked with
sour cream, urine, liquor, \gls{śuktā}, \gls{madhuka},
    \gls{sārivā}, and \gls{padmaka}.
    
\label{pradeha}
There are five types of \se{pradeha}{salve}, mashed with
\se{kṣārodaka}{alkaline water}, explained as follow: %
a paste of white \gls{sarṣapa}, %
a paste of sesame and \gls{aśvagandhā}, %
a paste of \gls{dāru} \gls{śelu}, and \gls{kapittha}, %
a paste of honey, \gls{śigru}, and \gls{punarnavā}, %
or with pastes of \gls{vyoṣa}, \gls{tiktā}, \gls{pṛthakparṇī}, and
\gls{bṛhatī}.\footnote{The Sanskrit syntax of this final poultice type does not
    exactly match with the previous items in the list.}
    
    
\paragraph{Congestion}
    
\item[11] 

And in the case of a \se{saṃsarga}{combination} or  a
\se{sannipāta}{congestion} of humours, one should use a mixture of the
ways of treatment that have been stated.\footnote{The distinction of
    \se{saṃsarga}{combination} and \se{sannipāta}{congestion} -- as two or
    three humours being reduced or irritated, respectively -- described in
    \CS\ \Ca{3.6.11}{254} and in \AH\ \Ah{1.1.12}{10}
    \parencite[tr.][206]{wuja-2003}. See also \volcite{1}[863,
    876]{josi-maha}.}

\item[12]

In all cases, one should consume \gls{harītakī} with jaggery.

Alternatively, one should drink \gls{pippalī}, ground in milk, with
five increments.\footnote{Ḍalhaṇa clarified that five extra
    \gls{pippalī} should be taken each day for ten days.} The food is milk
    and rice pudding for ten nights. Then one should reduce them by five
    incremental degrees.  In the same way until there are
    five.\footnote{The description of this procedure is not clear.} This
        is known as “The Long Pepper Incrementer”
        (\emph{Pippalīvarddhamānakam}). This cures wind-blood, irregular
        fever, loss of appetite, pallor, dropsy, piles, wheezing, cough,
        wasting disease, weak digestion, and heart disease.
        
The poultice is  \gls{sahā}, \gls{candana}, \gls{mūrvā}, \gls{piyāla},
\gls{śatāvarī}, \gls{kaśerukā}, \gls{sahadevā}, \gls{padmaka},
\gls{madhuka}, \gls{śatapuṣpā}, and \gls{kuṣṭha}. Alternatively, there
is a paste of \gls{śaileyaka}, \gls{aṭarūṣaka}, \gls{balā},
\gls{atibalā}, and \gls{jīvantī} mixed with milk.

Alternatively there is a paste of \gls{kāśmarī}, \gls{madhuka}, and \gls{tarpaṇa} mixed with the \se{maṇḍa}{residue} of ghee. 

Finally,\footnote{“Finally” expresses the final \dev{iti} in this
    sentence.  The vulgate omits this enclitic and adds \dev{sarveṣu ca}
    which \Dalhana{4.5.12}{426} interpreted as a statement that this
    massage therapy is useful for all five types of \se{vātarakta}{wind-blood}, whether caused by one to three humours.
    This idea does not seem to be present in the Nepalese version.}  there
    is the option of a \se{piṇḍataila}{sludgy oil} cooked with milk that
    is mixed with wax\sse{madhūcchiṣṭa}{wax}, \gls{śārivā},
    \gls{sarjarasa}, \gls{madhuka} and \gls{madhuravarga}.\footnote{The
        \CS\ discussed \dev{piṇḍataila} at \Ca{Ci.29.121--123}{632}, where
        Cakrapāṇidatta also described how to prepare it, with variant opinions
        from the authority Jatūkarṇa. \Dalhana{4.5.12}{426} said that it was
        an unfiltered oil with residue or, according to others, oil of
        \gls{śatāvarī} (\dev{piṇḍatailamiti aparīsrutaṃ kiṭṭayuktaṃ; kecit
        piṇḍataile śatāvarītyadhīyate/}). The vulgate makes it clearer that
        this is for \se{abhyaṅga}{oil massage}; see footnote
        \ref{massage-terms} for discussion.}

For the purpose of drinking, use old ghee that is cooked with
\gls{āmalaka} and \gls{sarala} and sweetened with sugar and honey.

For the purpose of \se{pariṣeka}{drizzling}, use old ghee that is cooked with
\gls{jīvanīya} or that is cooked with a decoction of
\gls{suṣavī}.\footnote{On \dev{jīvanīya} (\gls{jīvanīya}) see note
    \ref{jīvanīya} above. \Dalhana{4.5.12}{426} said that these compounds
    are kinds of ghee; this applies to the vulgate, not the Nepalese
    version.  \dev{pariṣeka}, also \dev{parīṣeka}, is described as a form of wash by \Dalhana{4.29.10--12}{504} and \Dalhana{6.39.156}{684}.  See \volcite{1}[479]{josi-maha}.}
    
            
Finally, cooked \gls{balā} oil is for drizzling, immersion, enema, and for
consumption.     
    
One should eat food preparations made of rice, \gls{ṣaṣṭika}, barley
and wheat accompanied with milk, meat broth or \gls{mudga} soup that
is not sour.   

And because of the diagnosis of an elevated humour, one should employ
blood-letting along with the therapies of vomiting, purging of bowels,
enema, and oily enema.\footnote{This is the “five therapy”
    (\dev{pañcakarma}) regime according to the \SS.  The five therapies
    were, in the tradition of the \CS, emetics, purgation, two types of
    enema, and nasal catharsis. The \SS\ replaced one of the enema
    treatments with bloodletting. Later authors introduced sweating and
    massage, as well as other therapies.
    
    The vulgate has an added verse at this point that gives a quick-and-simple cure for wind-blood.  The addition breaks 
    the continuity of the medical narrative; it was also not known to Ḍalhaṇa (\Su{4.5.13}{525}).}
 
 
     
\item[14]

There are verses on this.

\begin{sloka}
Physicians can easily cure recent-onset
wind-afflicted blood with compound treatments such as these.  
And long-standing ones can be managed.
\end{sloka}

\item[15--16]

\begin{sloka}    
    The following are recommended for wind-blood: oil-rubs, drizzles,  \sepl{pradeha}{salve}, the breezes of a
    fan,  delightful, large and draught-free shelters,  
        pillows that are soft on the cheek, 
        comfortable beds, 
       and gentle \se{saṃvāhana}{dry rubs}.\footnote{On
           “salves”, the types of salve were
           described above, p.\,\pageref{pradeha}.
           
           The expression \dev{vyajanānilāḥ} “breezes of a fan” appears to
           contradict the recommendation of a wind-free
           environment.}
\end{sloka}

\newpage

\item[17]

Finally,
\begin{sloka}
    One should avoid exercise, intercourse, anger, eating hot, sour,
or salty foods, sleeping during the day, and food that is slimy or
heavy.\footnote{\Dalhana{4.5.17}{427} noted that foods that are
    “\se{uṣṇa}{hot}” in this context are those with hot
    \se{vīrya}{potency}.}
\end{sloka}

\subsection{Seizures}

\item[18]

Now, one should treat a person who is affected with
\sepl{apatānaka}{seizure},\footnote{Cf.\ p.\,\pageref{apatānaka} and note
    \ref{apatānaka-note}.} as long as they do not have droopy eyes, crooked
    eyebrows, paralysed fingers, sweating, they are not trembling or babbling,
    and they do not fall out of bed and are not stretched
    backwards.\footnote{\label{bahirāyāmin}\Dalhana{4.5.18}{427} interpreted
        the term \dev{khaṭvāpātin} as “falling out of bed." “Stretched backwards”
        translates the  term \dev{bahirāyāmin}, which is more literally “stretched
        outwards."  The \CS\ described it as being stretched backwards like a bow
        (e.g.,  \Ca{6.28.46--48}{618}).  The condition is not described in detail
        in the main text of the \SS, though it may be the same as \dev{bāhyāyāma}
        (\SS\ \Su{2.1.57}{265}, tr.\,p.\,\pageref{bāhyāyāma} above.) %       
        % {Ḍalhaṇa interpreted it as \dev{paścānmukhagāmin}
        %        “walking while facing backwards.”  }
        It is commonly thought to refer to opisthotonos.}
%\begin{quote} % Dharmamitra tr. of 4.5.18
%    One should undertake the treatment of a patient suffering from apatānaka (convulsive seizure) provided they do not have drooping eyes, distorted eyebrows, stiffness of the penis, excessive perspiration, tremors, delirious speech, falling from the bed, or external extension of the limbs.
%    In that case, having first anointed the body with unctuous substances (sneha) and performed fomentation (sveda), one should treat the patient with a sharp nasal drop (avapīḍana) for the purpose of purifying the head. Subsequently, one should administer clear ghee (sarpiraccha) cooked with the decoction of the vidārigandhādi group, meat broth, milk, and curd; for in this way, the vāyu (vital wind) does not spread excessively. Then, having collected the bhadradārvādi group of vāta-subduing herbs, along with barley, kola, and kulattha, and the five groups of meat from marshy and aquatic animals, one should boil them together. Taking that decoction, one should mix it with sour liquids and milk, and cook it with ghee, oil, muscle fat (vasā), and marrow, adding the paste of the madhuraka group as an after-throw (pratīvāpa). This traivṛta (three-fold unctuous preparation) should be applied to patients with apatānaka for affusion (pariṣeka), immersion baths (avagāha), massage (abhyaṅga), drinking, food, non-unctuous enemas (āsthāpana), and nasal drops (nasya). One should also foment the patient according to the prescribed methods of fomentation. If the vāta is powerful, one should place the patient in a pit filled with comfortably warm chaff, husks, and dried cow-dung up to the mouth; or one should lay the patient upon a heated cart-furnace or a heated stone slab sprinkled with wine and covered with palāśa leaves; or one should foment the patient with kṛśara (rice and pulse porridge), āveśavāra (meat broth), and milk pudding (pāyasa).
%    Oil prepared with the fresh juices of radish (mūlaka), urubū, sphūrja, arjaka, arka, saptalā, and śaṅkhinī should be used for affusion and other treatments for those suffering from apatānaka. Sour curd mixed with black pepper and vacā, when drunk by one who has not eaten, cures apatānaka; or alternatively, oil, ghee, muscle fat, and honey may be used.
%    This protocol has been described for apatānaka caused by pure vāta; in cases of combined humors (saṃsṛṣṭa), a combined treatment should be administered. Between the paroxysms, one should administer nasal drops (avapīḍa). The patient should consume the muscle fat of roosters, crabs, black fish, porpoises, and boars, or milks processed with vāta-subduing drugs, or gruel (yavāgū) prepared with barley, kola, kulattha, radish, curd, ghee, and oil. Once the paroxysm has subsided after ten nights, one should treat the patient with unctuous purgatives (snehavirecana), non-unctuous enemas (āsthāpana), and unctuous enemas (anuvāsana). One should also observe the therapeutic measures prescribed for vātavyādhi (vāta diseases) and perform protective rituals (rakṣākarma). ||18||
%\end{quote}  
%\begin{quote} % Dharmamitra tr. of Ḍalhaṇa
%    He describes the treatment for apatānaka (convulsive seizure) starting with the words apatānakinam, etc.
%    Asrastākṣam: one whose eyes are not sunken or excessively displaced.
%    Apralāpinam: one who does not engage in incoherent or nonsensical speech.
%    Akhaṭvāpātinam: one who does not fall from the bed (khaṭvā). The meaning is this: one who, even while falling due to the onset of apatānaka, supports himself on the ground with both hands. One who habitually falls from the bed using both hands and feet as if they were the legs of the bed is a khaṭvāpātī; this patient is not of that type.
%    Abahirāyāminam: one who does not stretch out the limbs externally or move backward.
%    Avapīḍa: a preparation made by grinding head-purgative (śirovirecana) substances and squeezing the juice out; hence it is called avapīḍa. Even in apatānaka caused by pure vāta, head-purgation is administered here to immediately restore consciousness, as the vāta has reached the seat of kapha.
%    Anantaram: meaning "after" the head-purgation.
%    Regarding the ghṛta (clarified butter): it should be prepared with the four liquids—such as the decoction of the vidārigandhādi group—and with a paste (kalka) of the vidārigandhādi group added in a quarter-part (pādika).
%    Having described the ghṛta suitable for drinking, he speaks of the mahāsnēha (great oleation) useful for all procedures, starting with bhadra, etc.
%    Bhadradārvādi: the group beginning with bhadradāru, kuṣṭha, haridrā, varuṇa, etc.
%    Some authorities say that there should be four parts each of the decoction of the pañcavarga groups of bhadradārvādi, yavakōlādi, and marshy/aquatic animals (ānūpaudaka), as well as four parts each of kāñjika (sour gruel) and milk. Gayadāsa, however, considers the milk to be equal in quantity to the snēha (fat).
%    Jejjaṭa explains the "marshy/aquatic meat pañcavarga" as: those that frequent the banks (kūlacara), those that float (plava), those that live in shells (kośastha), those that walk with their feet (pādinaḥ), and fish. Gayadāsa, however, prefers the reading "along with the group of marshy/aquatic meats" (sānūpaudakamāṃsavargam), because by the term "marshy/aquatic" (sānūpaudaka), all five groups are included.
%    Amlāni: sour substances such as surā (wine), sauvīraka (a type of vinegar), dhānyāmla (fermented grain liquid), etc.
%    One should cook the sarpiḥ (clarified butter) together with oil, muscle fat (vasā), and marrow, adding an equal amount of each of the oil and other substances.
%    Madhurakaprativāpam: adding a paste of the kākolyādi group as a supplementary ingredient.
%    Traivr̥tam: it is called traivr̥ta because the clarified butter is "enveloped" (vr̥ta) by three other substances here.
%    Svēdavidhānaiḥ: by means of the methods of sudation described in the section on the application of sudation.
%    Nidadhyāt: one should place or apply it; this is a ūṣmasvēda (hot-bed sudation).
%    Rathakāraḥ: a blacksmith or metal-worker.
%    Kr̥śarā, etc.: this is an upanāhasvēda (poultice sudation).
%    Urubūka: white castor oil plant (śuklairaṇḍa). Sphūrjaka: shaped like a phaṇijjaka, generally found in the Pāriyātra mountains. Arjaka: kuṭhēraka. Śaṅkhinī: a variety of yavatiktā. Prepared with the juice of radish (mūlaka) and others in a quantity four times that of the snēha; having cooked the paste of radish and others in a quarter-part or without a paste, it is used lukewarm for sprinkling (pariṣeka) or liquid sudation (dravasvēda); by the word "etc." (ādi), it includes massage (abhyaṅga) and other uses.
%    Maricavacāyuktam: combined with black pepper and vacā; this refers to small amounts of vacā and pepper.
%    Saṃsṛṣṭe saṃsṛṣṭaṃ kartavyam: when there is a combination, a combined treatment should be applied. Upon perceiving a combination of the secondary factors (anubandha) of pitta and kapha, that very combined treatment which alleviates pitta and kapha should be administered. Jejjaṭa, however, says that in a combined state, treatment should be like that for vātarakta.
%    Vēgāntarēṣu: during the intervals between the paroxysms, when the paroxysm has subsided.
%    Tāmracūḍaḥ: a cock (kukkuṭa). Kr̥ṣṇamatsyaḥ: a black fish, scaleless and shaped like a crab.
%    The fat of the cock and others is for internal consumption. As stated by Vriddhavāgbhaṭa: "One should cause the patient to drink the fat of the cock, crab, porpoise (śiśumāra), and boar" (A.S. Ci. 23). ||18||
%\end{quote}  
%
In that context, at the very outset one should treat the oiled and
sweated patient with a sharp \se{avapīḍa}{sternutatory} in order to
clear the head.\footnote{The āyurvedic oil sternutatory is discussed in
    the \SS\  at \Su{4.40.44}{556} and its ingredients are listed at
    \Su{6.24.4}{650}.} %
    Next, the patient should be made to drink filtered ghee that is cooked
    with a decoction of herbs in the \gls{vidārigandhā} group,\footnote{For
        the plants in this group, see \cite[55]{suve-2000}.} sugarcane juice,
        milk, and yogurt. In that way, the wind does not spread excessively.
        
          % got to here
        \newpage
        \begin{tt}
            \raggedright
            \bigskip
            \marginpar{Draft tr.\ from here}
            
    
        
        
        

Thereafter, one should gather wind-alleviating herbs such as
\gls{bhadradāru}, etc. and other constituent parts, along with
\gls{yava}, \gls{kola}, and \gls{kulattha}, and the flesh of a
freshwater aquatic creature all at one place and prepare a decoction
of them. One should take this decoction and mix it properly with sour
substances and milk, and then cook the \textit{pratīvāpa}\footnote{It
    refers to an admixture of substances to medicines either during or
    after decoction. Refer to Monier-Williams's Sanskrit dictionary.} of
    \gls{madhuka} in this mixture along with ghee, oil, body fat, and bone
    marrow. This is \textit{trivṛt} that should be recommended in
    treatments of sprinkling, oil massage, applying a poultice, oral
    consumption, oily enema, and errhine for patients having spasmodic
    contractions.

The patient should then be made to sweat by the methods described
earlier. If the wind is stronger then the patient should be immersed
in [a vessel] filled with lukewarm fluid used for sprinkling
(\textit{trivṛt}). Or he should be kept in the hot fireplace of a
blacksmith.\footnote{H has the reading \dev{rathākāracullyām} that
    means \enquote{fireplace shaped like a chariot}, but the vulgate
    reading \dev{rathakāracullyām} makes more sense here. Thus, we have
    accepted it.} Or else he should be made to sweat by [a mixture of]
    \gls{kṛśara}, \textit{veśavāra},\footnote{Refer the above text no.7
        for \textit{veśavāra}. In H, the syllable \dev{vai} should have been
        \dev{ve}.} and milk.

Oil cooked with the juice of \gls{mūlaka}, \gls{uruvūka},
\gls{sphūrjaka}, \gls{saptalā}, and \gls{śaṅkhinī} should be used in
sprinking, etc. for patients with spasmodic contractions.\footnote{The
    word \dev{tailam} is not present in H but is present in the vulgate.
    We have accepted it.} Sour yogurt mixed with \gls{marica} and drunk on
    an empty stomach alleviates spasmodic contractions. Or else, ghee,
    oil, body fat, or bone marrow [can be consumed on an empty stomach].

This procedure of treatment thus described is for spasmodic
contractions caused only by wind. When mixed humours cause it then the
treatment should also be mixed. And when the spasms subside the
patient should be given \textit{avapīḍa}-s. One should also consider
the fats of cock, crab, black fish, and porpoise.\footnote{H has the
    reading \dev{rasān} which means \enquote{juices}. It seems unrealistic
    that juice would be extracted by crushing these whole animals. Vulgate
    has the reading \dev{vasāḥ} instead of \dev{rasān} which appears to be
    the more probable reading. Thus, we have accepted it.} Milk prepared
    with wind-alleviating medicines. Gruel prepared with barley,
    \gls{kola}, \gls{kulattha}, \gls{mūlaka}, yogurt, ghee, and oil.

One should treat this recurring spasm for ten nights with oil massage,
purging of bowels, enemas, and oily enemas. One should also look up
the treatment of diseases caused by wind. One should also undertake
preventive measures.

\item[19]

One should treat the paralytic (hemiplegic) patient whose limbs are
not languid, who is in pain, and who is self-composed. There, at the
beginning itself the patient should be massaged with oil and made to
sweat. After cleansing the patient with a mild
purifier,\footnote{According to P. V. Sharma, this refers to mild
    evacuatives (purgatives).} he should be administered with an oily
    enema and then a non-oily enema. Then at the appropriate time, he
    should be treated with special enemas of the brain and the head
    according to the method prescribed in the treatment of
    \textit{ākṣepaka}.\footnote{Refer \textit{Nidānasthāna} 1.50-51 for
        \textit{ākṣepaka}.}\q{Search for the section where the treatment of
            \textit{ākṣepaka} is described.} \textit{Anutaila} should be used for
        massage.\footnote{For the procedure of preparing \textit{anutaila},
            refer \textit{Cikitsāsthāna} 4.28.} \textit{Sālvala} should be used
            for poultice.\footnote{For the procedure of preparing
                \textit{sālvala}, refer \textit{Cikitsāsthāna} 4.14-15.} \gls{balā}
                oil should be used for oily enema.\q{Make the first letter of sentence
                    capital.} In this way, the unremitting patient should take the
                treatment for three to four months.

\item[20]
\item[21]

\subsubsection{Paralysis}

\item[22] One should treat the patient with \textit{ardita} %   
% \footnote{Refer
%    \textit{Nidānasthāna} 1.71-72 for
%\textit{ardita}.}
who is strong and possesses the means with the method prescribed in
treating wind diseases.\footnote{On \se{ardita}{paralysis}, see
    p.\,\pageref{ardita} above.} The unique thing is the treatment with
    enemas of the brain and the head, errhine, smoke, poultice, and steam
    bath through tubes. Then, one should take the great five roots
    (\textit{pañcamūlī}) with grass and prepare its decoction in milk
    mixed with twice the water. Then, the decoction with the milk
    remaining\footnote{It means that the water has evaporated.} should be
        brought down [the stove] and filtered. It should then be mixed with a
        \textit{prastha}\footnote{Ḍalhana commented \citep[425]{vulgate} that
            a \textit{prastha} is a measure of weight that is equal to 32
            \textit{pala-s}.} of oil and again placed over fire and cooked
            thoroughly. Then, the oil mixed with milk should be brought down [the
            stove] and then churned after it cools down. This is called
            \textit{kṣīrataila} that should be used in drinks, etc. for patients
            with \textit{ardita}.

\item[23]

In the diseases of \textit{gṛdhrasī}, \textit{viścañcī},
\textit{kroṣṭukaśīrṣa}, \textit{paṅgukalāya}, lameness,
\textit{vātakaṇṭaka}, burning sensation in the foot, numbness of the
foot, \textit{avabāhuka}, deafness, and \textit{dhamanīvāta}, one
should pierce the blood vessel as described earlier and, barring the
case of \textit{avabāhuka}, one should look up the treatment for wind
diseases.

\item[24]

However, in the case of \textit{karṇamūla},\footnote{The vulgate has
    the reading \dev{karṇaśūle} which appears to be a more credible
    reading according to the context.} lukewarm juice of
    \gls{śṛṅgavera}\footnote{\dev{śṛṅgavera} appears to be a name of
        ginger. Refer to the Sanskrit dictionary of Monier Williams.} mixed
        with \gls{madhuka}, oil, and salt should be put into the
        ears.\footnote{In H, the reading \dev{rasaiḥ} does not seem to make
            sense here. Hence we have accepted the vulgate reading \dev{rasam}.}
            Or else one can use goat urine, \gls{madhuka}, and oil. Or else one
            can use oil that is cooked with \gls{mātuluṅga}, \gls{dāḍima},
            \gls{tintiḍīka} juice, and urine.\footnote{In H, the word \dev{taila}
                should have been \dev{tailam} to make proper sense. The vulgate has
                this reading. Thus we have accepted it.} Or else one can use oil that
                is cooked with sour liquor, buttermilk, and urine.

One should also make the patient sweat with a steam bath through
tubes. One should also look up the treatment for wind diseases. More
will be said later.

\item[25]

In the case of \textit{tūnī} and \textit{pratitūnī}, one should make
the patient drink ghee and salt with hot water. Or else one should
administer the powder of \gls{pippalī} and other herbs with hot water.
Or else one should make the patient drink ghee that is made thick with
asafoetida and \gls{yavakṣāra}.\footnote{\dev{yavakṣāra} is an alkali
    prepared from the ashes of burnt green barleycorns. Refer to the
    Sanskrit dictionary of Monier Williams.} One should also treat the
    patient with enemas.

\item[26]

In the case of \textit{ādhmāna},\footnote{Refer to
    \textit{Nidānasthāna} 1.88. V. S. Apte explains it as
    \enquote{swelling of the belly}. P.V. Sharma has translated it as
    flatulence.} however, one should do \textit{avatarpaṇa},\footnote{We
        are unclear about its meaning. The vulgate has the reading
        \dev{apatarpaṇa} that means fasting.} heating the hands,
        \textit{phalavartikriyā},\footnote{The entry \dev{phalavarti} has the
            meaning \enquote{suppository} in the Sanskrit dictionary of Monier
            Williams. The Cambridge dictionary explains suppository as \enquote{a
            small, solid pill containing a drug that is put inside the anus, where
            it dissolves easily.} Refer to the link
            https://dictionary.cambridge.org/dictionary/english/suppository. Last
            accessed 30-Oct-2023.} stimulation of digestion, and [administer]
            digestives. One should also employ the purging of bowels and enemas.

In the case of \textit{pratyādhmāna},\footnote{Refer to
    \textit{Nidānasthāna} 1.89. According to the Sanskrit dictionary of
    Monier Williams, it is a kind of tympanites or wind-dropsy.} one
    should employ vomiting, fasting, and stimulation of digestion.

\item[27]

In the case of \textit{aṣṭhīlā} and
\textit{pratyaṣṭhīlā},\footnote{Refer to \textit{Nidānasthāna} 1.90
    and 1.91.} the procedure is that of \textit{gulma} and internal
    abscess.

\item[28]

The beneficial asafoetida, the three pungent spices (long pepper,
black pepper, and dry ginger), \gls{vacā}, \gls{ajamoda} grains,
\gls{ajagandhā}, \gls{dāḍima}, \gls{tintiḍīka}, \gls{pāṭhā},
\gls{citraka}, \gls{cavyā}, \gls{saindhava}, \gls{viḍa},
\gls{sauvarcala}, \gls{yavakṣāra}, \gls{suvarcikā}, \gls{pippalī}
root, \gls{amlavetasa}, \gls{śaṭhī}, \gls{puṣkara}, \gls{hapuṣa}
(juniper berry), and \gls{ajājī} (cumin seeds) should be powdered.
This powder should be mixed with a lot of \gls{mātuluṅga} juice. Then
it should be made into pills each weighing one \textit{akṣa}.
Thereafter the patient of wind disease should consume one pill every
morning. This medicine indeed cures \textit{gulma}, rapid breathing,
cough, loss of appetite, heart disease, \textit{ādhmāna},
\textit{pārśvodara}, \textit{bastiśūla}, \textit{anāhamūtra}, painful
piles, \textit{plīhodara}, and \textit{pāṇḍuroga}. Also, this medicine
is excessively used in cases of \textit{tūnī} and \textit{pratitūnī}.

\item[29]

There are verses in this regard.

\begin{sloka}
The wind that has entered into the body tissues should be correctly
understood as either pure or vitiated by humours\footnote{In H, the
    reading \dev{lakṣaṇonyām̐śca} does not make sense. Hence I cannot
    translate it. Perhaps the correct reading could be
    \dev{lakṣaṇābhyāñca}. This would connect with the two conditions of
    the wind as stated in the verse.} and should be cured accordingly.
\end{sloka}

\item[30]

\begin{sloka}
The wind that is accompanied by fat causes a swelling that is painful,
hard, and cold. The physician should properly treat it like a treating
a swelling.
\end{sloka}

\item[31]

\begin{sloka}
When the wind accompanied by phlegm and fat enters the thighs, it
causes pain in and immobility of the thighs due to numbness, pain, and
fever.
\end{sloka}

\item[32]

\begin{sloka}
Also, the thighs become pained, stiff, cold, and do not quiver due to
sleep. They become heavy and as if belonging to someone
else.\footnote{In H, the verb \dev{vartate} should have been in the
    dual. Also, the word \dev{āsthirau} does not make sense. The vulgate
    has the sensible reading \dev{asthirau} which we have accepted here.}
\end{sloka}

\item[33]

\begin{sloka}
That is called \textit{ūrūstambha}. Others call it \textit{āḍhyavāta}.
In that case, one should drink the \textit{ṣaṇḍharaṇa} powder with
cool water.
\end{sloka}

\item[34]

\begin{sloka}
Similarly, consuming the powder of \gls{pippalī} and other herbs with
hot water is beneficial. Or else, one should consume the powder of
\textit{triphalā} with \gls{kṣaudra} and \gls{kaṭukā}.
\end{sloka}

\item[35-38]

\begin{sloka}
Or else, one should drink the best \gls{guggulu} or \gls{śilājatu}
with urine. Such a person cures the wind that is afflicted by phlegm
and accompanied by fat, as well as heart disease, loss of appetite,
\textit{gulma}, and internal abscess.

One should employ salty urine [therapy], sudation, and hard rubbing.
One should also apply [the paste of ] mustard and \gls{karañja} fruits
mixed with urine.\footnote{The word \dev{dihet} in H is not a proper
    Sanskrit word. We have taken its proper form \dev{dihyāt} as given in
    the vulgate.}

One should eat old \gls{śyāmāka}s, \gls{kodrava}, \gls{uddālaka}, etc.
along with uncooked\footnote{The vulgate has the reading
    \dev{aghṛtaiḥ} that means without ghee.} flesh of wild animals and
    unsalted vegetables that are beneficial.
\end{sloka}

\item[39]

\begin{sloka}
When the phlegm and fat become amply reduced one should again employ
the treatment of oil massage, etc. for the patient.
\end{sloka}
\end{tt}

\end{translation}
